\chapter{Background and Literature Review}
\label{chap:background}

\section{Sports Analytics and Biomechanical Feedback}
\label{sec:sports_analytics}

The application of quantitative analysis to athletic performance has grown substantially over the past two decades, driven by the availability of sensor data, video, and machine learning. In basketball, analytics has historically focused on game outcomes---shot selection efficiency, player tracking, and play-by-play statistical modelling---rather than on the mechanical correctness of individual skills~\cite{Goldman2013}. Biomechanical analysis, by contrast, operates at the sub-second level, examining joint angles, velocity profiles, and positional trajectories during a single skill execution.

Research by Cabarkapa et al.~\cite{Cabarkapa2021} established normative joint angle ranges for proficient and non-proficient free throw shooters, providing the first rigorously validated dataset linking specific angular measurements to shooting performance. Their work identified elbow extension at release (optimal 170--180\textdegree), elbow flexion in the preparatory phase (optimal 75--85\textdegree), and forearm verticality (optimal 0--8\textdegree) as the most discriminative features between high and low accuracy groups.

Okazaki et al.~\cite{Okazaki2012,Okazaki2015} analysed the effect of shooting distance on optimal release angle, establishing that the ideal arc angle decreases from approximately 78\textdegree{} at close range (2.8\,m) to approximately 66\textdegree{} at mid-range (4.6\,m) and 58\textdegree{} at long range (6.0\,m+), reflecting the trade-off between arc height and horizontal velocity. This distance-dependent parameterisation is critical for any system that aims to assess outdoor or variable-range shooting.

Struzik et al.~\cite{Struzik2014} contributed a biomechanical analysis of the jump shot, documenting the angular change in wrist position between release and follow-through completion (optimal 18--28\textdegree) as a measurable proxy for the wrist snap action that experienced coaches identify qualitatively.

\section{Computer Vision in Sports}
\label{sec:cv_sports}

Computer vision has been applied to sports analysis along two main axes: player tracking and action recognition. Player tracking systems---most prominently the optical tracking systems deployed in NBA arenas since 2013---capture multi-player positional data at high frequency but provide no information about joint-level mechanics. Action recognition approaches classify observed motion into discrete categories (e.g., ``made shot,'' ``dribble,'' ``pass'') but do not recover the continuous joint angle trajectories needed for biomechanical assessment~\cite{Rao2020}.

The specific problem of shot mechanics analysis requires a different approach: recovering per-frame joint positions with sufficient accuracy and temporal density to compute angle time series, then detecting the physiologically meaningful events (crouch, release, landing) within that series.

\section{Pose Estimation Approaches}
\label{sec:pose_estimation}

Human pose estimation has evolved through several architectural generations, each improving accuracy and inference speed:

\textbf{OpenPose}~\cite{Cao2017} introduced bottom-up part affinity field (PAF) estimation, allowing multi-person pose detection from a single forward pass. While accurate, OpenPose requires GPU inference and does not produce reliable depth estimates from monocular images.

\textbf{HRNet}~\cite{Sun2019} introduced high-resolution representation learning for human pose estimation, maintaining spatial resolution across all network stages rather than using a typical encoder--decoder structure. HRNet achieves state-of-the-art keypoint detection accuracy on COCO benchmarks but requires significant computational resources.

\textbf{MediaPipe Pose}~\cite{Lugaresi2019}, developed by Google, takes a different trade-off: it sacrifices some accuracy relative to HRNet in order to run in real time on CPU hardware. The model produces 33 full-body landmarks with associated visibility confidence scores, which is sufficient for the joint subset required by basketball shooting analysis. Critically, it operates on standard RGB input, requires no calibration, and runs at practical speed on the same server hardware that serves the API. For these reasons, SHOOTRZ uses MediaPipe 0.10.14 as the primary pose estimator, with HRNet present in the repository as a candidate for a future upgrade.

\section{Object Detection and Ball Tracking}
\label{sec:object_detection}

YOLOv8~\cite{Jocher2023}, developed by Ultralytics, is the eighth generation of the YOLO (You Only Look Once) real-time object detection family. The nano variant (YOLOv8n, 6.5\,MB) achieves fast inference with a modest accuracy reduction compared to larger model sizes, making it suitable for CPU deployment. SHOOTRZ uses YOLOv8n to optionally detect the basketball across sampled frames and fit a parabolic trajectory for release angle estimation. Ball detection is treated as a supplementary signal because it requires the ball to be visible and clearly distinguishable in the frame---a constraint that fails frequently in amateur recording conditions.

\section{Signal Processing for Biomechanical Time Series}
\label{sec:signal_processing}

Raw pose keypoints extracted frame-by-frame contain measurement noise arising from MediaPipe's probabilistic inference, JPEG compression artefacts, and motion blur. Direct differentiation of noisy position signals to obtain velocities or angle derivatives amplifies this noise significantly. The Savitzky--Golay filter~\cite{Savitzky1964} addresses this by fitting a polynomial of specified degree to a sliding window of data points using a least-squares criterion, preserving the shape of peaks and troughs while attenuating high-frequency noise. SHOOTRZ applies a Savitzky--Golay filter with window length 5 and polynomial order 2 to each joint coordinate series independently, chosen to preserve the sharp wrist-height peak at release while suppressing inter-frame jitter.

\section{Large Language Models in Domain-Specific Applications}
\label{sec:llm}

Large language models (LLMs) have demonstrated an ability to generate domain-specific guidance when prompted with structured numerical context. In the sports domain, early work has explored LLM-generated player performance summaries and strategy briefs, though the use of LLMs to produce structured biomechanical coaching feedback from computed metrics is, to the authors' knowledge, novel at the time of writing.

\begin{sloppypar}
A key challenge in applying LLMs to this domain is the tendency to generate plausible but vague responses when numerical specificity is required. SHOOTRZ addresses this by constraining Gemini 2.5 Flash~\cite{Gemini2024} to fill a defined Pydantic output schema---\lstinline|ShotFeedbackOutput|---rather than generating free-form text. The schema includes specific fields for metric-level verdicts, ranked strengths, ranked improvements, and a numerical score, ensuring that the LLM's output is directly consumable by the mobile frontend without further parsing.
\end{sloppypar}

\section{Recommendation Systems in Sports}
\label{sec:recommendation}

Drill and exercise recommendation in sports coaching is a multi-objective problem: recommendations must be relevant to the player's identified weaknesses, appropriately challenging for their skill level, and varied enough to prevent repetitive training. Traditional rule-based systems (``if elbow angle is low, recommend elbow extension drill'') are brittle and cannot adapt based on which drill types have historically produced improvement for similar users.

\begin{sloppypar}
FAISS (Facebook AI Similarity Search)~\cite{Johnson2019} provides efficient nearest-neighbour search over dense vector embeddings, enabling retrieval of the most mechanically similar drills to a user's current profile. The contextual bandit framework---specifically the LinUCB algorithm~\cite{Li2010}---extends retrieval-based recommendation with a principled exploration--exploitation trade-off, selecting the drill cluster that maximises expected reward (skill improvement) conditioned on the user's mechanical feature vector. This combination, implemented via the MABWiser library, is the core of SHOOTRZ's recommendation engine.
\end{sloppypar}

\section{Existing Systems and Related Work}
\label{sec:related_work}

Several commercial and research systems address aspects of the basketball shooting analysis problem:

\textbf{HomeCourt} (Apple App Store) uses on-device computer vision to count makes and misses and measure shot arc in real time. It does not compute joint angles or provide biomechanical metric breakdowns.

\textbf{ShotTracker} is an arena-level system using wristband sensors and basket-mounted hardware to track shot attempts and outcomes. It requires dedicated hardware and is not applicable to individual practice sessions.

\textbf{Coach's Eye} is a video annotation tool that allows coaches to draw overlays and record voiceover commentary on uploaded video. It does not perform automated pose estimation or metric computation.

\textbf{Kinovea} is an open-source sports video analysis tool that supports manual angle measurement from video but provides no automated pose estimation or scoring.

None of the above systems combines automated biomechanical metric extraction, research-validated scoring, natural-language coaching generated by a large language model, and a mobile-accessible interface. SHOOTRZ is designed to close this gap specifically for basketball shooting mechanics.

\section{Summary}
\label{sec:background_summary}

This chapter has established the research foundations and the technology landscape from which SHOOTRZ is built. The biomechanics literature provides well-validated normative ranges for the key joint angles in the shooting motion. MediaPipe Pose provides a practical foundation for monocular pose estimation on commodity hardware. LLMs provide structured natural-language feedback generation when properly constrained with output schemas. The FAISS and LinUCB approaches provide a principled foundation for personalised drill recommendation. Chapter~\ref{chap:srs} now specifies the full system requirements derived from these foundations.
